\documentclass[12pt, a4paper]{article}

\usepackage[utf8]{inputenc}
\usepackage[T1]{fontenc}
\usepackage{amsmath, amssymb, amsthm}
\usepackage{geometry}
\usepackage{fancyhdr}
\usepackage{enumitem}
\usepackage{xcolor}
\usepackage{hyperref}

\geometry{margin=1in}
\pagestyle{fancy}
\fancyhf{}
\rhead{线性代数与常微分方程}
\lhead{JingXueQing}
\rfoot{Page \thepage}

\newcommand{\problem}[1]{\subsection*{Problem #1}}
\newcommand{\solution}{\paragraph{Solution.}}

\definecolor{solutioncolor}{RGB}{0, 100, 0}

\begin{document}

\begin{center}
    {\LARGE\bfseries 习题 7 解答} \\[0.5em]
    {\large 线性代数与常微分方程} \\[0.3em]
    JingXueQing \quad | \quad 2026-10-07
\end{center}

\hrule
\vspace{1em}
\problem{1}

设矩阵

$$A = \begin{bmatrix} 4 & 1 \\ 2 & 3 \end{bmatrix}$$

求 $A$ 的行列式 $\det(A)$ 与秩 $\operatorname{rank}(A)$，并判断 $A$ 是否可逆。

---

\solution

用按 $ad-bc$ 二阶公式计算 $\det(A) = 10$。

对 $A$ 求行最简形，非零行数为 $2$，故 $\operatorname{rank}(A) = 2$。

结果已通过上式直接计算得到，可回代原式核对。

\textbf{Answer:}
\[\boxed{\det(A) = 10;\quad \operatorname{rank}(A) = 2}\]

\vspace{1em}
\hrule
\vspace{1em}

\problem{2}

设

$$A = \begin{bmatrix} 2 & 1 \\ 1 & 2 \end{bmatrix}$$

求 $A$ 的全部特征值与对应特征向量，并判断 $A$ 是否可以对角化。

---

\solution

构造特征矩阵 $A - \lambda I$。

$\det(A - \lambda I) = \left(\lambda - 3\right) \left(\lambda - 1\right)$。

令其为零解得特征值：$\lambda = 3, 1$

求每个特征值对应的特征向量：$\lambda=1$ 时特征向量张成 $\mathrm{span}\{\left[\begin{matrix}-1\\1\end{matrix}\right]\}$（维数 1）；$\lambda=3$ 时特征向量张成 $\mathrm{span}\{\left[\begin{matrix}1\\1\end{matrix}\right]\}$（维数 1）

把各特征向量按列排成 $P = \left[\begin{matrix}-1 & 1\\1 & 1\end{matrix}\right]$。

则 $D = P^{-1}AP = \left[\begin{matrix}1 & 0\\0 & 3\end{matrix}\right]$，即 $A = PDP^{-1}$。

\textbf{Answer:}
\[\boxed{P^{-1}AP = D = \left[\begin{matrix}1 & 0\\0 & 3\end{matrix}\right]}\]

\vspace{1em}
\hrule
\vspace{1em}

\problem{3}

求矩阵

$$B = \begin{bmatrix} 1 & 2 \\ 2 & 1 \end{bmatrix}$$

的逆矩阵 $B^{-1}$，并验证 $BB^{-1} = I$。

---

\solution

对增广矩阵 $(A\mid I)$ 施以初等行变换，把左半部分化为单位阵，右半部分即为 $A^{-1}$。

$A^{-1} = \left[\begin{matrix}- \frac{1}{3} & \frac{2}{3}\\\frac{2}{3} & - \frac{1}{3}\end{matrix}\right]$。

回代验证：$AA^{-1} = \left[\begin{matrix}1 & 0\\0 & 1\end{matrix}\right]$，等于单位阵。

\textbf{Answer:}
\[\boxed{\left[\begin{matrix}- \frac{1}{3} & \frac{2}{3}\\\frac{2}{3} & - \frac{1}{3}\end{matrix}\right]}\]

\vspace{1em}
\hrule
\vspace{1em}

\problem{4}

判断向量组

$$\mathbf{v}_1 = \begin{bmatrix} 1 \\ 2 \\ 3 \end{bmatrix}, \quad
  \mathbf{v}_2 = \begin{bmatrix} 2 \\ 4 \\ 6 \end{bmatrix}, \quad
  \mathbf{v}_3 = \begin{bmatrix} 1 \\ 0 \\ 1 \end{bmatrix}$$

是否线性无关，并说明理由。

---

\solution

以各向量为列构成 $3\times3$ 矩阵（共 3 个向量）。

计算得 $\operatorname{rank}(A) = 2$。

秩 $2$ < 向量个数 $3$ ⟹ 线性相关（存在非平凡线性组合为零）　（$\det A = 0$）。

例如存在非平凡线性组合 - 2\,\left[\begin{matrix}1\\2\\3\end{matrix}\right] + 1\,\left[\begin{matrix}2\\4\\6\end{matrix}\right] + 0\,\left[\begin{matrix}1\\0\\1\end{matrix}\right] = 0，故线性相关。

\textbf{Answer:}
\[\boxed{\text{线性相关}}\]

\vspace{1em}
\hrule
\vspace{1em}

\problem{5}

解一阶常微分方程

$$y' + 2y = 0$$

求其通解。

---

\solution

识别方程（移项后）：$2 y{\left(x \right)} + \frac{d}{d x} y{\left(x \right)} = 0$。

调用符号求解器 dsolve 求通解。

通解：$y{\left(x \right)} = C_{1} e^{- 2 x}$。

题目未给初始条件，故给出通解。

**残差验证**：将解代回原方程，残差恒为 0，解正确。

\textbf{Answer:}
\[\boxed{y{\left(x \right)} = C_{1} e^{- 2 x}}\]

\vspace{1em}
\hrule
\vspace{1em}

\problem{6}

解二阶常系数齐次线性微分方程

$$y'' + y = 0$$

求其通解。

---

\solution

识别方程（移项后）：$y{\left(x \right)} + \frac{d^{2}}{d x^{2}} y{\left(x \right)} = 0$。

调用符号求解器 dsolve 求通解。

通解：$y{\left(x \right)} = C_{1} \sin{\left(x \right)} + C_{2} \cos{\left(x \right)}$。

题目未给初始条件，故给出通解。

**残差验证**：将解代回原方程，残差恒为 0，解正确。

\textbf{Answer:}
\[\boxed{y{\left(x \right)} = C_{1} \sin{\left(x \right)} + C_{2} \cos{\left(x \right)}}\]

\vspace{1em}
\hrule
\vspace{1em}

\problem{7}

一质量为 $2\,\mathrm{kg}$ 的物体受到合力 $F = 10\,\mathrm{N}$ 的作用，
求该物体的加速度 $a$（用牛顿第二定律）。

---

\solution

牛顿第二定律 $F = ma$。

由题面读得 $F = 10\,\mathrm{N}$，$m = 2\,\mathrm{kg}$。

$a = \frac{F}{m} = \frac{10}{2} = 5\,\mathrm{m/s^2}$。

\textbf{Answer:}
\[\boxed{a = 5\,\mathrm{m/s^2}}\]

\vspace{1em}
\hrule
\vspace{1em}

\problem{8}

电路中电压 $U = 220\,\mathrm{V}$ 加在电阻 $R = 55\,\mathrm{ohm}$ 两端，
求通过电阻的电流 $I$（用欧姆定律）。

---

\solution

欧姆定律 $I = \dfrac{V}{R}$。

$V = 220\,\mathrm{V}$，$R = 55\,\mathrm{\Omega}$。

$I = \dfrac{220}{55} = 4\,\mathrm{A}$。

\textbf{Answer:}
\[\boxed{I = 4\,\mathrm{A}}\]

\vspace{1em}
\hrule
\vspace{1em}

\problem{9}

两个点电荷 $q_1 = 2\,\mathrm{C}$ 与 $q_2 = 3\,\mathrm{C}$ 相距 $r = 1\,\mathrm{m}$，
求它们之间的库仑力大小。

---

\solution

库仑定律 $F = k\dfrac{q_1q_2}{r^2}$，$k = 8.99\times10^9\,\mathrm{N\,m^2/C^2}$。

$q_1 = 2\,\mathrm{C}$，$q_2 = 3\,\mathrm{C}$，$r = 1\,\mathrm{m}$。

$F = 8.99\times10^9 \times \dfrac{2\times3}{1^2} = 53940000000\,\mathrm{N}$。

\textbf{Answer:}
\[\boxed{F = 53940000000\,\mathrm{N}}\]

\vspace{1em}
\hrule
\vspace{1em}

\problem{10}

一个物体从 $h = 45\,\mathrm{m}$ 高处自由落下，
求它落地所需的时间与落地速度（取 $g = 9.8\,\mathrm{m/s^2}$）。

---

\solution

取竖直向下为正方向，重力加速度 $g = 9.8\,\mathrm{m/s^2}$（取自题面）

由自由落体公式 $h = \tfrac12 g t^2$ 解出下落时间。

$t = \sqrt{\frac{2h}{g}} = \sqrt{\frac{2 \times 45}{9.8}} = 3.03\,\mathrm{s}$。

落地速度 $v = gt = 9.8 \times 3.03 = 29.7\,\mathrm{m/s}$。

\textbf{Answer:}
\[\boxed{t = 3.03\,\mathrm{s},\quad v = 29.7\,\mathrm{m/s}}\]

\vspace{1em}
\hrule
\vspace{1em}

\problem{11}

（超出当前求解器支持范围，用于测试诚实失败行为）
请证明：对任意 $n \times n$ 可逆矩阵 $A$，其特征值均不为零。

---

\textit{\color{red} 未求解: LLM 已禁用，且该题不属确定性题库}

\vspace{1em}
\hrule
\vspace{1em}

\problem{12}

设某区间上的分段函数为

$$f(x) = \begin{cases} x^2 & x < 0 \\ -x & x \ge 0 \end{cases}$$

求 $\lim_{x \to 0} f(x)$，并判断 $f$ 在 $x = 0$ 处是否连续。

\textit{\color{red} 未求解: LLM 已禁用，且该题不属确定性题库}

\vspace{1em}
\hrule
\vspace{1em}


\vspace{2em}
\noindent\textit{Auto-generated: 10/12 problems solved by SymPy.}
\end{document}
